\section{Procedure}\label{sec:procedure}


%Presents the apparatus and procedure used and includes:
%1) a list of the apparatus;
%2) an apparatus diagram; 3) the technical details of the investigation in the procedure used to acquire data.
%Written in the past tense with specific terminology, this should allow others to replicate the experiment.


The experimental apparatus and procedure followed the standard method outlined in the University of Victoria PHYS 326
laboratory manual~\cite{phys326}.

Before any data was collected, the diameter of the toroid was measured using a 30cm ruler.
The diameter was determined from measuring the middle-middle of the shape in order to derive the total length, $L$.
The Hall sensor, which measures the magnetic field intensity, was positioned between the gap of the toroid seen in
Figure~\ref{fig:toroid}, and angled perpendicular to the incoming field lines.

The LoggerPro software was then sut up to measure the voltage across the coil as well as manual entries from the
Gaussmeter.
The Gaussmeter measurements were converted to units of Teslas.

A calculated column was used to determine $H$ using Eq.~\eqref{eq:H}.
A technique was employed in an attempt to demagnetise the metal sample.
First, the coil was energised up to 2.0 Amps to substantially polarise the material, then switched to a reverse
current of the same magnitude using the commutator switch to reverse the polarity.
This was repeated in 0.1 Amp increments all the way down to zero amps.
After four attempts of this technique, the lowest reading from the Gaussmeter was 11.3 G.

Starting from zero Amps, the data was collected in small intervals by entering in the digital voltage and Gaussmeter
readings into LoggerPro.
Once a maximum current of 2.0 Amps was reached, data was taken as current was reduced incrementally back down to zero.
At that point, the commutator switch was flipped to reverse the direction of current, and the process was repeated
to capture more data points to a maximum of 2.0 Amps.
Current was again reduced to zero where the data approximately closed the hysteresis curve.

\begin{table}[htbp]
    \caption{Experimental Parameters}
    \begin{tabular}{lrrr}
        \toprule
        {Quantity}           & {Reading}                \\
        \midrule
        $r$, Toroid radius   & \qty{82(5)}{\milli\meter} \\
        $N$, Number of turns & 1000                     \\
        $R$, Coil resistance & \qty{3.488(0.010)}{\ohm} \\
        \bottomrule
    \end{tabular}\label{tab:params}
\end{table}

Table~\ref{tab:params} provides a list of measurements taken.

The measurement uncertainty of $\pm 5$ millimeters was used to account for the uneven shape of the electrical tape
around the toroid and the ruler millimeter graduations.
The measurement uncertainty of resistance is taken from the DMM's spec sheet.